\chapter{Introduction}\doublespacing
\label{Chapter1}
\lhead{Chapter I.}

Credit risk prediction is one of the oldest and most economically
consequential applications of statistical learning. Whenever a lender
chooses whether to extend credit to an applicant, the decision is
backed --- explicitly or implicitly --- by a model that assigns each
applicant a probability of default. A classifier that wrongly clears
a defaulter incurs a direct monetary loss, while one that wrongly
rejects a creditworthy applicant loses future revenue and damages the
lender's competitiveness. The two errors are not symmetric in cost,
yet both must be controlled at the same time.
Because the consequences of a misclassification fall on real money,
the development of credit-risk classifiers has long served as a
proving ground for new statistical and machine-learning ideas.

The literature on credit scoring covers a wide spectrum of methods:
classical linear discriminant analysis and logistic regression, non-parametric methods
such as $k$-nearest neighbours, decision trees and rule-based systems, and modern ensembles and neural networks
\citep{breiman2001random,chen2016xgboost}. Across this
body of work two observations recur:
\begin{enumerate}
 \item No single algorithm is uniformly best across all credit
 datasets; the winner depends on dataset size, class balance,
 feature mix and the metric of interest.
 \item Combining several heterogeneous classifiers --- through bagging,
 boosting or stacked generalisation --- tends to produce more
 robust performance than any single learner alone
 \citep{wolpert1992stacked}.
\end{enumerate}
Both observations motivate the present work. We do not propose a
single new classifier; rather, we build a \emph{stacked ensemble} of
existing classifiers, select its second-level (meta) learner by an
objective rule, and submit the entire benchmark to non-parametric
significance testing across multiple datasets so that the empirical
ordering of the classifiers is supported by formal statistics rather
than by raw accuracy figures alone.

\section{Background and Motivation}
\label{sec:motivation}

A credit-risk classifier maps an applicant's covariate vector
$\mathbf{x}\in\mathbb{R}^{p}$ to a binary label
$y\in\{0,1\}$ (default or no default). In practice the classifier
produces a score $s(\mathbf{x})\in[0,1]$ interpreted as the estimated
probability of default; a threshold $t$ is then applied to obtain a
crisp class label. Although this formulation is simple, three issues
make credit datasets challenging in practice:

\textbf{Class imbalance.} Defaulters are typically the minority class.
A naive classifier that predicts the majority label for every
applicant achieves high accuracy but is useless for the lender, who
specifically needs to identify the minority of likely defaulters.
The standard responses are either to resample the training data ---
for example by the \emph{SMOTE} oversampling scheme of
\citet{chawla2002smote} --- or to tune the decision threshold $t$
explicitly, as in the \emph{GHOST} procedure of
\citet{esposito2021ghost}.

\textbf{Mixed feature types.} Credit applications contain numeric
attributes (age, income, debt ratio) alongside categorical ones
(occupation, marital status, residence type). A robust pipeline
needs both an encoding scheme that respects this distinction and a
scaling step that prevents the numeric attributes from dominating
distance-based classifiers such as the $k$-nearest-neighbours rule.

\textbf{Threshold and metric sensitivity.} Several widely-reported
metrics --- accuracy, $F_1$, Cohen's $\kappa$ --- depend on a chosen
threshold $t$. Two classifiers can produce identical area under the
ROC curve (AUC), which is threshold-independent, and yet report very
different accuracies because their default thresholds happen to favour
different operating points. This is a known source of confusion in
the benchmark literature, and our evaluation
explicitly reports both threshold-dependent and threshold-independent
metrics to guard against it.

A further complication, often glossed over in single-dataset studies,
is statistical: when several classifiers are compared on several
datasets, the differences in observed performance must be assessed
relative to the variability across datasets, not in absolute terms.
\citet{demvsar2006statistical} laid out the now-standard
non-parametric procedure for this purpose: a Friedman rank test
followed by an appropriate post-hoc correction --- typically Holm
\citep{holm1979simple} or Nemenyi --- to pin down which pairs of
classifiers are significantly different. The present thesis applies
this protocol to a four-dataset benchmark.

\section{Problem Statement and Objectives}
\label{sec:problem}

The empirical question of this thesis is straightforward:
\begin{quote}
\emph{Among a representative set of credit-risk classifiers, which
model performs most consistently across multiple available UCI
datasets, and is its advantage statistically significant relative to
the others?}
\end{quote}
The answer is operationalised by the following objectives:

\begin{itemize}
 \item[\textbf{1.}] Implement eight representative classifiers ---
 Logistic Regression, $k$-Nearest Neighbours, Support Vector Machine
 (RBF kernel), Decision Tree, Random Forest, Gradient Boosting,
 XGBoost, and a Stacked Ensemble --- and apply them uniformly to
 four UCI credit datasets (Australian, Japanese, German, Taiwan).
 \item[\textbf{2.}] Build a feature-selection sweep with four
 strategies (Full features, Recursive Feature Elimination,
 L1-penalised Logistic Regression, Forward Selection) per dataset,
 so that the comparison spans both ``unselected'' and ``selected''
 feature regimes.
 \item[\textbf{3.}] Make the Stacked Ensemble's second-level learner
 a choice from five candidates (Logistic Regression, Random Forest,
 XGBoost, Multi-Layer Perceptron, and a Genetic-Algorithm-tuned
 logistic regression, GA-MaxAccLogit), and select the meta by an
 objective $F_1$-macro rule on $k$-fold out-of-fold predictions.
 The GA-MaxAccLogit meta is adapted from the ProfLogit method of
 \citet{stripling2018proflogit}: we keep ProfLogit's genetic
 machinery but replace its expected-profit objective with $F_1$-macro.
 \item[\textbf{4.}] Tune Random Forest and XGBoost by \emph{Bayesian
 optimisation} (Gaussian-process surrogate, Expected-Improvement
 acquisition) so that the heavyweight learners are not unfairly
 penalised by default settings \citep{snoek2012practical}.
 \item[\textbf{5.}] Apply \emph{SMOTE} to the training data of every
 imbalanced dataset and tune the decision threshold via \emph{GHOST}
 on the same datasets, to guarantee that the minority class is not
 silently ignored by the threshold-dependent metrics.
 \item[\textbf{6.}] Evaluate every (dataset, feature-selection,
 classifier) triple on the held-out test set using F1-macro, AUC and
 a ten-measure diagnostic suite, and report the
 full leaderboard.
 \item[\textbf{7.}] Test the global null hypothesis ``all eight
 classifiers perform equally'' by the \emph{Friedman} rank test
 \citep{friedman1937use} across the sixteen evaluation blocks
 (4 datasets $\times$ 4 feature-selection variants), and resolve
 the pairwise comparisons by the \emph{Wilcoxon signed-rank} test
 with \emph{Holm} step-down correction \citep{demvsar2006statistical}.
\end{itemize}

Objectives \textbf{1}--\textbf{6} construct the benchmark;
objective \textbf{7} elevates the comparison from descriptive to
inferential. The combination is what distinguishes the present
thesis from many credit-scoring benchmarks in the literature.

\section{Contributions}
\label{sec:contributions}

The principal contributions of this thesis are the following:

\begin{enumerate}
 \item A \emph{uniform pipeline} that applies the same preprocessing,
 feature-selection sweep, hyperparameter tuning, and evaluation
 suite to all eight classifiers on all four UCI datasets, producing
 a leaderboard of $4\times4\times8=128$ model rows that is
 internally consistent.
 \item An adaptation of the ProfLogit genetic-algorithm logistic
 model \citep{stripling2018proflogit}, here called
 \textbf{GA-MaxAccLogit}, in which the expected-profit fitness
 function is replaced by the classification metric $F_1$-macro
 (penalised by an $L_1$ norm on the coefficients). The adapted
 learner is plugged in as a candidate meta-learner of the stacked
 ensemble.
 \item A statistical-significance layer using \emph{Friedman} +
 \emph{Wilcoxon-Holm} on the rank matrix of the sixteen
 (dataset $\times$ FS) blocks, treating each block as an
 independent evaluation. This is the methodology of
 \citet{demvsar2006statistical} applied at the
 block --- not dataset --- level, and it yields markedly more
 power than the equivalent four-dataset test.
 \item A reproducible \emph{Colab notebook bundle} containing the
 Friedman test, the pairwise Wilcoxon-Holm post-hoc, the
 GA-MaxAccLogit implementation, and the data-loading scripts
 --- all runnable end-to-end on the open UCI data.
\end{enumerate}

\section{Organisation of the Thesis}
\label{sec:organisation}

The remainder of this thesis is organised as follows.

\textbf{Chapter \ref{Chapter2}} reviews the eight classifiers used
in the benchmark and develops the mathematical formulation of each.
The stacked ensemble is introduced at the end of the chapter as the
formal meta-learner setup that ties the previous seven together.

\textbf{Chapter \ref{Chapter3}} describes the methodological
contributions: the four feature-selection strategies, the SMOTE +
GHOST imbalance pipeline, the Bayesian-optimisation tuning protocol,
and the GA-MaxAccLogit meta-learner adapted from ProfLogit.

\textbf{Chapter \ref{Chapter4}} describes the four UCI datasets,
the preprocessing pipeline, the evaluation metrics, and the
implementation environment.

\textbf{Chapter \ref{Chapter5}} reports the empirical results: the
per-dataset leaderboards, the cross-dataset comparison, the Friedman
test, and the Wilcoxon-Holm pairwise post-hoc. The chapter ends with
a discussion of why the Stacked Ensemble is the most consistent
performer.

\textbf{Chapter \ref{Chapter6}} concludes the thesis, discusses the
limitations of the present study, and outlines directions for future
work. Appendix \ref{AppendixA} reproduces the principal Python code
listings used in the experiments.
